%% kapa: exam=1 round=26 year=2015 period=2 form=A subject=appraisal_law subject_ko=부동산관계법규 seq=40 number=40
%% source: Q-Net 감정평가사 기출문제 (cst00309) artlSeq=5138001 file=extracted/1008/5138001/제26회_감정평가사 1차_기출/제26회_감정평가사 1차_기출/2교시_감평_A형.pdf
%% conversion: kordoc 4.24.51 -> tools/gen1.py (auto); review=proofread-sonnet; flags=none
\begin{question}{40}
부동산등기법령상 담보권 등기에 관한 설명으로 옳은 것은?
\begin{choices}
  \item 등기원인에 그 약정이 있는 경우에도 변제기는 저당권의 등기사항이 아니다.
  \item 변제기와 이자의 약정이 있는 경우에도 그 내용은 저당권부채권에 대한 질권의 등기사항이 아니다.
  \item 등기관이 일정한 금액을 목적으로 하지 아니하는 채권을 담보하기 위한 저당권설정의 등기를 할 때에는 그 채권의 평가액을 기록하지 아니한다.
  \item 등기관이 동일한 채권에 관하여 여러 개의 부동산에 관한 권리를 목적으로 하는 저당권설정의 등기를 할 때에는 각 부동산의 등기기록에 그 부동산에 관한 권리가 다른 부동산에 관한 권리와 함께 저당권의 목적으로 제공된 뜻을 기록하여야 한다.
  \item 등기관이 채권의 일부에 대한 양도로 인한 저당권 일부이전등기를 할 때에는 권리에 관한 등기사항이외에 별도로 양도액을 기록하지 아니한다.
\end{choices}
\end{question}
\answer{4}
